\section{Definitive teachings that are not proposed as revealed dogma}

Some truths are so connected with Revelation that the deposit cannot be
faithfully guarded or expounded without them. When the Church proposes such a
truth definitively, it is to be held firmly and irrevocably. The assent differs
formally from divine and Catholic faith because the proposition is not proposed
as itself formally revealed; its certainty is nevertheless definitive.

The 1998 commentary gives examples while warning readers not to confuse
examples with an exhaustive catalog. The category includes truths connected
logically or historically with Revelation and judgments required to preserve
the deposit. The same doctrine may later be proposed as revealed if the Church
authoritatively determines that formal relation; classifications follow the
Church's act, not a writer's preference.

The distinction refutes a common binary: ``infallible dogma or optional
opinion.'' It also restrains maximalism. A Catholic cannot infer definitiveness
from importance alone. The competent authority must manifest the definitive
proposal, whether by a defining act or by an authoritative confirmation that
the ordinary and universal Magisterium has already taught the matter
infallibly.

